%====================================
% KAPITEL 2: ARBEITSTITEL, FORSCHUNGSFRAGEN UND ZIEL
%====================================
\section{Arbeitstitel}

Der Arbeitstitel dieser Bachelorarbeit lautet: \glqq Welche Change-Management-Faktoren beeinflussen den Erfolg von ERP-Implementierungen in mittelst\"andischen Unternehmen?\grqq{} Er fasst die drei Kernelemente zusammen: Change Management als Untersuchungsgegenstand, ERP-Implementierungen als technologischen Kontext und den Mittelstand als spezifisches Umfeld. Gepr\"ufte Alternativen -- eine deskriptive Kurzform oder ein Fokus auf Einzelfaktoren -- bilden diese Breite nicht ab.

\section{Forschungsfragen und Ziel der Arbeit}

Es folgt die Ableitung der zentralen Forschungsfragen sowie die Definition des \"ubergeordneten Ziels. Ausgehend von der identifizierten Forschungsl\"ucke im Bereich der ERP-Implementierungen im deutschen Mittelstand wird zun\"achst die Hauptforschungsfrage formuliert. Anschlie\ss{}end werden vier Sub-Fragen entwickelt, die verschiedene Analyseperspektiven auf das Thema er\"offnen.

\subsection{Hauptforschungsfrage}

Zentral ist die Beantwortung der Frage:

\begin{quote}
\textbf{Welche Change-Management-Faktoren beeinflussen den Erfolg von ERP-Implementierungen in mittelst\"andischen Unternehmen?}
\end{quote}

Sie adressiert die Schnittstelle zwischen zwei etablierten Forschungsfeldern, die bislang nur selten gemeinsam f\"ur den spezifischen Kontext mittelst\"andischer Unternehmen betrachtet wurden. W\"ahrend sowohl die ERP-Erfolgsforschung als auch das Change-Management f\"ur sich genommen umfangreich untersucht sind, fehlt es an integrierenden Modellen, die beide Perspektiven systematisch verbinden. Insbesondere die Frage, welche spezifischen Change-Management-Faktoren unter den besonderen Rahmenbedingungen mittelst\"andischer Unternehmen -- flache Hierarchien, knappe Ressourcen, geringere IT-Affinit\"at -- tats\"achlich erfolgskritisch sind, ist bislang nicht hinreichend beantwortet.

\subsection{Sub-Fragen}

Aus der Hauptforschungsfrage leiten sich vier Sub-Fragen ab, die jeweils einen spezifischen Teilaspekt des Untersuchungsgegenstands beleuchten.

\textbf{Sub-Frage 1:} Welche Dimensionen des DeLone \& McLean IS-Success-Modells sind f\"ur den Kontext mittelst\"andischer Unternehmen besonders relevant und wie l\"asst sich der Erfolg einer ERP-Implementierung anhand dieses Modells operationalisieren?

Das DeLone \& McLean-Modell\footcite[Vgl.][S.~9--30]{delone2003} gilt als eines der am besten validierten Erfolgsma\ss{}e im Bereich der Informationssysteme. Allerdings wurde es \"uberwiegend in gro\ss{}en Unternehmen und internationalen Kontexten getestet. Die Anwendung auf ERP-Systeme im Mittelstand erfordert daher eine kritische Pr\"ufung der Dimensionen und gegebenenfalls eine Anpassung an die Spezifika dieser Zielgruppe. Die relative Gewichtung der einzelnen Dimensionen (Systemqualit\"at, Informationsqualit\"at, Servicequalit\"at, Nutzerzufriedenheit, Nettonutzen) unterscheidet sich dabei mutma\ss{}lich von den in der Gro\ss{}unternehmensforschung dokumentierten Auspr\"agungen und bedarf einer gesonderten Betrachtung.

\textbf{Sub-Frage 2:} Welche Change-Management-bezogenen Erfolgsfaktoren f\"ur ERP-Implementierungen lassen sich aus der aktuellen Forschungsliteratur ab 2015 identifizieren?

Der zeitliche Rahmen ab 2015 ist bewusst gew\"ahlt, da sich die Anforderungen an ERP-Projekte durch Digitalisierung, Cloud-L\"osungen und agile Methoden in den letzten Jahren grundlegend ver\"andert haben. \"Altere Studien sind oft noch von traditionellen, monolithischen Einf\"uhrungsprojekten gepr\"agt, die nur noch bedingt auf die heutige Praxis \"ubertragbar sind. Durch eine systematische Literaturanalyse sollen die zentralen Faktoren extrahiert, kategorisiert und hinsichtlich ihrer empirischen Fundierung bewertet werden.

\textbf{Sub-Frage 3:} Welche Rolle spielen Kommunikation, F\"uhrungsunterst\"utzung und Schulungsma\ss{}nahmen im Rahmen von ERP-Implementierungen in mittelst\"andischen Unternehmen und wie lassen sich diese Ma\ss{}nahmen anhand des ADKAR-Modells bewerten?

Diese drei Faktoren gelten in der Literatur als besonders erfolgskritisch. Allerdings fehlen bislang detaillierte Untersuchungen dazu, wie sie speziell im Mittelstand wirken und welche Auspr\"agungen sie annehmen m\"ussen, um nachhaltig zum Projekterfolg beizutragen. Das ADKAR-Modell\footcite[Vgl.][S.~124--132]{galli2018} bietet hierf\"ur einen strukturierten Bewertungsrahmen, der die f\"unf Phasen des individuellen Ver\"anderungsprozesses abbildet und so eine differenzierte Analyse erm\"oglicht.

\textbf{Sub-Frage 4:} Welche konkreten Handlungsempfehlungen f\"ur die Planung und Durchf\"uhrung von Change-Management-Ma\ss{}nahmen in ERP-Einf\"uhrungsprojekten mittelst\"andischer Unternehmen lassen sich aus den Befunden der Literaturanalyse ableiten?

Die Ergebnisse sollen nicht nur wissenschaftlichen Erkenntnisgewinn liefern, sondern auch praktisch nutzbar sein. Die abgeleiteten Empfehlungen sollen Entscheidungstr\"agern in mittelst\"andischen Unternehmen eine Orientierung bieten und dazu beitragen, die hohen Ausfallraten von ERP-Projekten zu reduzieren.

\subsection{Ziel der Arbeit}

Ziel der Untersuchung ist es, einen literaturbasierten Beitrag zum Verst\"andnis der Erfolgsfaktoren von ERP-Implementierungen im deutschen Mittelstand zu leisten. Im Kern soll ein konzeptioneller Bezugsrahmen entwickelt werden, der die Erfolgsdimensionen nach DeLone \& McLean mit den Change-Management-Faktoren aus der aktuellen Forschung und dem ADKAR-Modell verbindet. Auf dieser Basis sollen praxisrelevante Handlungsempfehlungen abgeleitet werden, die Unternehmen dabei unterst\"utzen, ihre ERP-Einf\"uhrungsprojekte durch gezielte Change-Management-Ma\ss{}nahmen erfolgreicher zu gestalten.
